\documentclass[11pt]{article}
\usepackage[margin=1in]{geometry}
\usepackage{amsmath,amssymb}
\usepackage{graphicx}
\usepackage[hidelinks]{hyperref}

\title{Literature-Implied Answer: The \(p=\infty\) Endpoint of Reinov's Proposition 3.6}
\author{Automatic Math Discovery Run \texttt{fa\_banach\_001}}
\date{June 18, 2026}

\begin{document}
\maketitle

\section*{Source Question}

The source paper is O. I. Reinov, \emph{Approximation of operators in Banach
spaces}, arXiv:1002.3902 \cite{Reinov2010}. Proposition 3.6 states that for
every \(p\in[1,\infty)\) there are separable reflexive Banach spaces \(X\)
and \(Y\) such that
\[
\overline{X^*\otimes Y}^{\tau_p}=X^*\widehat{\widehat\otimes}_pY.
\]
Remark 3.7 then says that it is seemingly unknown whether this remains true
for \(p=+\infty\). In the parsed source this is
\path{data/parsed/arxiv_sources/1002.3902/source.tex}, lines 1365--1390.

\begin{center}
\includegraphics[width=0.95\linewidth]{figures/open_problem_crop.png}
\end{center}

\section*{Identification}

The endpoint follows from Pitt's theorem. Choose real numbers
\[
1<s<r<\infty,\qquad X=\ell_r,\qquad Y=\ell_s.
\]
Then \(X\) and \(Y\) are separable and reflexive. For \(p=\infty\), the
absolutely \(\infty\)-summing norm is the operator norm, so
\[
\Pi_\infty(X,Y)=\mathcal L(X,Y).
\]
Moreover, \(\tau_\infty\) is the topology of compact-open convergence.

Finite-rank operators are compact-open dense in \(\mathcal L(X,Y)\): given
\(T\in\mathcal L(X,Y)\), a compact set \(K\subset X\), and \(\varepsilon>0\),
choose a finite-rank operator \(P:X\to X\) with
\(\sup_{x\in K}\|Px-x\|<\varepsilon/\|T\|\), and then \(TP\) is finite-rank
and approximates \(T\) uniformly on \(K\).

By Pitt's theorem, every bounded operator \(\ell_r\to\ell_s\), \(1<s<r<\infty\),
is compact \cite{Pitt1936}. Since \(\ell_s\) has the approximation property,
every compact operator from \(\ell_r\) into \(\ell_s\) is a norm limit of
finite-rank operators. Hence
\[
\mathcal L(\ell_r,\ell_s)
=\mathcal K(\ell_r,\ell_s)
=\overline{\ell_r^*\otimes\ell_s}^{\,\|\cdot\|}.
\]
Combining this with compact-open density gives
\[
\overline{X^*\otimes Y}^{\tau_\infty}
=\mathcal L(X,Y)
=X^*\widehat{\widehat\otimes}_\infty Y.
\]
Thus Proposition 3.6 is true at \(p=+\infty\) for the pair
\((X,Y)=(\ell_r,\ell_s)\), \(1<s<r<\infty\).

\section*{Status and Scope}

Status: \texttt{literature\_implied\_answer}. This answers the existence
question in Reinov's Remark 3.7. It is not presented as a new theorem beyond
the standard Pitt theorem plus the endpoint interpretation of
\(\Pi_\infty\) and \(\tau_\infty\). The supporting literature did not
explicitly identify itself as answering Reinov's remark; the implication is
agent-identified.

\section*{Search Evidence}

The run indexes were searched for arXiv:1002.3902, approximation of operators,
operator approximation, compact approximation, bounded approximation property,
local reflexivity, Pitt, and the \(p=\infty\) endpoint terminology. No prior
packet or attempt for this source question was found.

\begin{thebibliography}{9}
\bibitem{Reinov2010}
O. I. Reinov,
\emph{Approximation of operators in Banach spaces}, arXiv:1002.3902, 2010.

\bibitem{Pitt1936}
H. R. Pitt,
\emph{A note on bilinear forms}, J. London Math. Soc. 11 (1936), 174--180.
\end{thebibliography}

\end{document}
